\begin{textsample}{POS Dim 1 – 2023\_c5 – Score 22.00 – 2023\_c5/enade\_2023\_arqui\_urb\_f2c5983230.txt}  \label{ex:f1_pos_008}
CAPÍTULO 4 DISTRIBUIÇÃO DE CONCEITOS A \textbf{finalidade} deste capítulo é apresentar um \textbf{panorama} nacional sobre a distribuição do Conceito Enade 2023.
Para traçar esse \textbf{panorama}, recorre -se a tabelas e gráficos, que são acompanhados de comentários e análises que contribuem para a contextualização dos dados neles expostos.
Para \textbf{visualizar} o \textbf{panorama} do Enade 2023, em nível nacional, na primeira seção deste capítulo, os dados sobre o número de cursos participantes são organizados por grande região.
Na \textbf{segunda}, o mapeamento do conceito do exame é feito por categoria administrativa das IES e modalidade de oferta dos cursos, de maneira \textbf{cruzada} às grandes regiões.
Na última seção, a lógica de apresentação dos dados é feita por organização acadêmica e, também, por grande região de funcionamento dos cursos, no país como um todo. 4.1 \textbf{PANORAMA} NACIONAL DA DISTRIBUIÇÃO DOS CONCEITOS A lógica adotada na apresentação dos dados, nesta seção e nas duas seguintes, pressupõe o registro de duas informações de natureza \textbf{reiterativa}.
O primeiro diz respeito aos valores possíveis do Enade, que variam em níveis de 1 a 5, sendo que, à medida que esses valores aumentam, o curso é melhor avaliado, e ao fato de que aqueles cursos que tiveram somente um ou nenhum concluinte participante não obtêm o Conceito Enade, ficando sem conceito ( SC ), como ressaltado na apresentação deste relatório.
O segundo registro é que, caso haja diferença entre os cursos \textbf{tabulados} neste capítulo e no capítulo 1, ela corresponderá aos cursos que não foram avaliados, ou seja, àqueles, em princípio, sem estudantes concluintes que tivessem sido inscritos para a prova.
Na Tabela 4.1, são apresentadas a quantidade e a distribuição de cursos de Arquitetura e Urbanismo participantes do Enade 2023, por faixa de conceito e grande região.
A partir dos dados da Tabela 4.1, nota -se que, dos 564 cursos participantes, 259 ( 45,9 \% ) classificaram -se com conceito 3, o valor modal.
Esse foi também o conceito modal em todas as regiões : Centro-Oeste ( 44,7 \% ), Nordeste ( 44,2 \% ), Norte ( 47,5 \% ), Sudeste ( 45,1 \% ), Sul ( 48,5 \% ).
O conceito 2 foi o segundo mais frequente em nível nacional ( 25,9 \%, correspondendo a 146 cursos ), tendo sido o segundo conceito mais frequente na maioria das regiões, exceto no Sul.
O conceito 4 foi o terceiro ( 17,2 \%, correspondendo a 97 cursos ).
Houve, ainda, 36 ( 6,4 \% ) cursos que receberam o conceito 5 e 13 ( 2,3 \% ) cursos que receberam o conceito 1.
Dos 564 cursos de Arquitetura e Urbanismo, 13 ( 2,3 \% ) ficaram sem conceito ( SC ).
A região Norte participou com 40 cursos, o que corresponde a 7,1 \% do total nacional.
Desses, 19 cursos receberam o conceito 3, o que \textbf{equivale} a 47,5 \% do total regional, o conceito modal na região, conforme mencionado.
A 14 cursos ( 35,0 \% ) atribuiu -se o conceito 2 ; a quatro cursos ( 10,0 \% ), o conceito 4 ; e a dois cursos ( 5,0 \% ) o conceito 1.
Ainda nessa região, nenhum curso ( 0,0 \% ) obteve conceito 5 e um curso ( 2,5 \% ) ficou sem conceito ( SC ).
A região Nordeste participou com 95 cursos, correspondentes a 16,8 \% do total nacional.
Nessa região, 42 cursos ( 44,2 \% em termos regionais ) obtiveram conceito 3, o conceito modal para a região, como já comentado.
Os conceitos 2 e 4 foram atribuídos, respectivamente, a 25 e a 15 cursos ( 26,3 \% e 15,8 \%, respectivamente ).
Ainda nessa região, sete cursos ( 7,4 \% ) foram avaliados com o conceito 5 e três cursos ( 3,2 \% ) foram avaliados com o conceito 1.
Outros três cursos ( 3,2 \% ) ficaram sem conceito ( SC ).
Para os 246 cursos participantes da região Sudeste ( 43,6 \% do total nacional ), o conceito 3 foi o conceito modal, como já anteriormente assinalado, tendo sido obtido por 111 cursos ( 45,1 \% ).
O conceito 2 foi atribuído a 67 cursos ( 27,2 \% ) e o conceito 4, a 44 cursos ( 17,9 \% ).
Receberam os conceitos 5 e 1, respectivamente, 15 cursos ( 6,1 \% ) e quatro cursos ( 1,6 \% ).
Nessa região, cinco cursos ( 2,0 \% ) ficaram sem conceito ( SC ).
Os 136 cursos da região Sul corresponderam a 24,1 \% do total nacional.
Como já \textbf{indicado}, houve predominância do conceito 3 ( conceito modal ), atribuído a 48,5 \% dos cursos da região, ou seja, atribuído a 66 dos 136 cursos participantes na região Sul.
O conceito 4 foi atribuído a 30 cursos ( 22,1 \% ) ; o conceito 2, a 26 cursos ( 19,1 \% ) ; e o conceito 5, a nove cursos ( 6,6 \% ).
Na região Sul, dois cursos ( 1,5 \% ) obtiveram conceito 1 e três cursos ( 2,2 \% ) ficaram sem conceito ( SC ).
Já dos 47 cursos participantes na região Centro-Oeste ( 8,3 \% do total nacional ), 21 ( 44,7 \% em termos regionais ) receberam conceito 3, o conceito modal, como já destacado. 14 cursos ( 29,8 \% ) obtiveram o conceito 2, cinco cursos ( 10,6 \% ) obtiveram o conceito 5 e quatro cursos ( 8,5 \% ), o conceito 4.
Nessa região, dois cursos ( 4,3 \% ) obtiveram o conceito 1 e um curso ( 2,1 \% ) ficou sem conceito ( SC ).
Em um gráfico de distribuição \textbf{cumulativa}, em cada valor de \textbf{abscissa} ( eixo x ), o valor apresentado no eixo das \textbf{ordenadas} ( eixo y ) é o valor acumulado da \textbf{variável} ( Conceito Enade ).
Em linhas gerais, em um gráfico desta mesma \textbf{tipologia}, \textbf{poligonais} mais à esquerda representam, no caso dos dados apresentados nesta seção, as grandes regiões com a pior distribuição, e \textbf{poligonais} mais à direita, as grandes regiões com os melhores conceitos.
A partir destas diretrizes, no Gráfico 4.1, é apresentada a distribuição \textbf{cumulativa} do Conceito Enade, segundo a grande região.
No Gráfico 4.1, observando a linha vertical \textbf{pontilhada}, que ajuda a \textbf{visualizar} o percentual de conceitos acumulados até o valor 3, ou seja, conceitos 1, 2 e 3, pode -se notar, por exemplo, que a região Norte ( linha azul ) apresenta 89,7 \% dos cursos com conceito 3 ou abaixo deste valor.
A região Sul ( linha \textbf{roxa} ), por outro lado, apresenta 70,7 \% dos cursos na mesma situação : com conceito 3 ou abaixo.
Destaca -se que, na região Norte, o conceito 5 não foi alcançado por nenhum curso, diferentemente das regiões Centro-Oeste ( linha laranja ), Nordeste ( linha verde ), Sudeste ( linha vermelha ) e Sul ( linha \textbf{roxa} ), nas quais todos os cursos têm conceito menor ou igual a 5, por isso as respectivas \textbf{poligonais} \textbf{cumulativas} dessas regiões alcançam 100 \% no conceito 5.
A região Norte ( linha azul ) apresenta os cursos com conceitos mais baixos, ao passo que a região Sul ( linha \textbf{roxa} ) apresenta a distribuição com valores maiores.
Os cursos em 4.2 CONCEITOS POR CATEGORIA ADMINISTRATIVA, MODALIDADE DE OFERTA DOS CURSOS E GRANDE REGIÃO Na Tabela 4.2, é apresentada a distribuição dos cursos de Arquitetura e Urbanismo participantes do Enade 2023, por categoria administrativa e modalidade de oferta, de acordo com os conceitos por eles alcançados, segundo as grandes regiões brasileiras.
De acordo com as informações da Tabela 4.2, dos 564 cursos participantes, 69 ( 12,2 \% ) eram \textbf{ministrados} em instituições públicas e 495 ( 87,8 \% ), em instituições privadas.
Quanto à modalidade de oferta, 544 ( 96,5 \% ) cursos eram presenciais e 20 ( 3,5 \% ), a distância.
De acordo com as informações da mesma tabela, em termos nacionais, entre instituições públicas e privadas, observa -se que, dos 36 cursos avaliados com conceito 5, 31 foram oferecidos em IES públicas, e 5, em IES privadas.
Dos 495 cursos participantes de IES privadas, o conceito 3 foi o valor modal, atribuído a 244 ( 49,3 \% ) cursos.
Em relação aos demais cursos de instituições privadas, 144 cursos ( 29,1 \% ) foram avaliados com o conceito 2, 76 cursos ( 15,4 \% ) foram avaliados com o conceito 4, 13 cursos ( 2,6 \% ) foram avaliados com o conceito 1 e, cinco cursos ( 1,0 \% ), com o conceito 5.
Nessa categoria, 13 cursos ( 2,6 \% ) ficaram sem conceito ( SC ).
Na rede pública, o conceito modal foi o 5, atribuído a 31 cursos, correspondentes a 44,9 \% dos 69 cursos da categoria.
Dos demais cursos participantes, 21 cursos ( 30,4 \% ) obtiveram o conceito 4, 15 cursos ( 21,7 \% ) receberam o conceito 3 e dois cursos ( 2,9 \% ), o conceito 2.
O conceito 1 foi atribuído a nenhum curso ( 0,0 \% ), conforme mencionado.
Nessa categoria administrativa, nenhum curso ( 0,0 \% ) ficou sem conceito ( SC ).
Dos 544 cursos participantes da modalidade presencial, o conceito 3 foi o valor modal, atribuído a 252 ( 46,3 \% ) cursos.
Dos demais cursos dessa modalidade de oferta, 134 cursos ( 24,6 \% ) receberam o conceito 2 ; 96 cursos ( 17,6 \% ) receberam o conceito 4 ; 36 cursos ( 6,6 \% ), o conceito 5 ; e 13 cursos ( 2,4 \% ), o conceito 1.
Nessa modalidade de oferta, 13 cursos ( 2,4 \% ) ficaram sem conceito ( SC ).
Dos 20 cursos participantes da modalidade a distância, o conceito 2 foi o valor modal, atribuído a 12 ( 60,0 \% ) cursos.
Dos demais cursos \textbf{ofertados} nessa modalidade, sete cursos ( 35,0 \% ) receberam conceito 3 e um curso ( 5,0 \% ) recebeu o conceito 4.
Nenhum curso ( 0,0 \% ) recebeu o conceito 5, o conceito 1 ou ficou sem conceito ( SC ).
Na análise por região, observa -se, que, na região Norte, as instituições públicas participaram com sete cursos ( 17,5 \% total regional ).
O conceito modal foi 3, atribuído a quatro cursos, correspondendo a 57,1 \% dessa categoria na região.
O conceito 4 teve uma concentração muito próxima do modal, com três cursos ( 42,9 \% ).
Nenhum curso recebeu os conceitos 5, 2 ou 1.
As instituições privadas participaram com 33 cursos na região Norte ( 82,5 \% do total regional ), e um curso ficou sem conceito ( SC ).
O conceito modal nas instituições privadas na região foi o 3, atribuído a 15 cursos, o que corresponde a 45,5 \% dessa categoria na região.
Os conceitos 2, 1 e 4 foram atribuídos a 14, dois e um curso, respectivamente.
Nenhum curso recebeu o conceito 5 ou ficou sem conceito ( SC ).
Dos cursos oferecidos na região Norte, 37 eram presenciais e um ficou sem conceito ( SC ).
Dos 37 cursos oferecidos, 19 obtiveram o conceito 3 ( modal ) ; 11, o conceito 2 ; quatro, o conceito 4 ; dois, o conceito 1 ; e nenhum, o conceito 5.
Nenhum curso ficou sem conceito ( SC ) na região Norte e na modalidade presencial.
Três cursos eram na modalidade a distância, tendo três recebido o conceito modal 2.
Nenhum curso ficou sem conceito ( SC ) ou recebeu os demais conceitos na região Norte e na modalidade a distância.
Na região Nordeste, a rede privada concentra 80 dos 95 cursos participantes, o equivalente a 84,2 \% do total da região.
Entre os cursos oferecidos por IES privadas no Nordeste, 39 cursos foram avaliados com conceito 3, o conceito modal. 25 cursos obtiveram o conceito 2 ; nove cursos, o conceito 4 ; três cursos, o conceito 1 ; e um curso obteveeram o conceito 5.
Nessa categoria administrativa e grande região, três cursos ficaram sem conceito ( SC ).
As instituições públicas da região Nordeste participaram com 15 cursos ( 15,8 \% do total da região ).
Desses, seis ( 40,0 \% ) obtiveram conceito 5, o valor modal, seguidos por seis cursos, que receberam o conceito 4.
O conceito 3 foi atribuído a três cursos ; o conceito 2, a nenhum curso ; e o conceito 1, a nenhum curso.
Nenhum curso ficou sem conceito ( SC ).
Dos cursos oferecidos na região Nordeste, 94 eram presenciais : três ficaram sem conceito ( SC ), três obtiveram o conceito 1 ; 25, o conceito 2 ; 41, o conceito 3, o valor modal ; 15, o conceito 4 ; e sete, o conceito 5.
Um curso era na modalidade a distância, tendo recebido o conceito modal Nenhum curso ficou sem conceito ( SC ) ou recebeu os demais conceitos na região Nordeste e na modalidade a distância.
Na região Sudeste, a proporção de cursos da rede privada, 91,1 \%, correspondeu a 224 dos 246 cursos participantes.
Nessa combinação de categoria administrativa e grande região, o conceito modal foi o 3, atribuído a 108 cursos ; e cinco cursos ficaram sem conceito ( SC ).
Os demais foram avaliados com os seguintes conceitos : 2 ( 66 cursos ), 4 ( 39 cursos ), 1 ( quatro cursos ) e 5 ( dois cursos ).
Dos 22 cursos oferecidos por instituições públicas na região Sudeste, 13 obtiveram o conceito 5 ; cinco, o conceito 4 ; três, o conceito 3 ; um obteve conceito 2 e nenhum obteve o conceito 1.
Nenhum curso ficou sem conceito ( SC ).
Ainda na região Sudeste, grande parcela dos cursos ( 97,6 \% ) foi oferecida na modalidade presencial : 240 dos 246 cursos.
Para esses, a moda foi o conceito 3, com 109 cursos.
Os demais cursos presenciais da região foram alocados nos conceitos 2 ( 64 cursos ), 4 ( 43 cursos ), 5 ( 15 cursos ) e 1 ( quatro cursos ).
Cinco cursos presenciais ficaram sem conceito ( SC ).
Na região Sudeste, seis cursos foram oferecidos na modalidade a distância, tendo sido nenhum com conceito 1 ; três com conceito 2 ; dois com conceito 3 ; um com conceito 4 ; e nenhum com conceito 5.
Nenhum curso ficou sem conceito ( SC ).
As instituições privadas concentraram 120 dos 136 cursos participantes da região Sul, 88,2 \% do total regional.
Desses, 63 cursos obtiveram conceito 3, o conceito modal.
Nessa combinação de categoria administrativa e grande região, 25 receberam o conceito 4 ; 25 receberam o conceito 2 ; dois, o conceito 5 ; e dois obtiveram o conceito 1.
Três cursos ficaram sem conceito ( SC ).
As instituições públicas na região Sul participaram com 16 cursos ( 11,8 \% dos cursos da região ), aos quais foram atribuídos os conceitos 1 ( nenhum curso ), 2 ( um curso ), 3 ( três cursos ), 4 ( cinco cursos ) e 5 ( sete cursos ).
Nenhum curso da região Sul oferecido por IES públicas ficou sem conceito ( SC ).
Na região Sul, 127 dos cursos foram oferecidos na modalidade presencial.
Para esses, a moda foi o conceito 3, com 62 cursos.
Os demais cursos presenciais da região foram alocados nos conceitos 4 ( 30 cursos ), 2 ( 21 cursos ), 5 ( nove cursos ) e 1 ( dois cursos ).
Ficaram sem conceito ( SC ) três cursos.
Na modalidade a distância, foram oferecidos nove cursos, sendo nenhum curso com conceito 1, cinco cursos com conceito 2, quatro cursos com conceito 3 ( modal ), nenhum curso com conceito 4 e nenhum curso com conceito 5.
Nenhum curso ficou sem conceito ( SC ) nessa região.
Na região Centro-Oeste, 38 dos 47 cursos participantes eram de IES privadas ( 80,9 \% em termos regionais ).
Desses, 19 cursos receberam conceito 3, o valor modal.
Dos demais cursos, 14 receberam o conceito 2, dois receberam o conceito 4, dois receberam o conceito 1 e nenhum recebeu o conceito 5.
Um curso ficou sem conceito ( SC ).
Quanto aos nove cursos oferecidos por instituições públicas na região Centro-Oeste ( 19,1 \% do total regional ), nenhum ficou com conceito 1 ; nenhum, com conceito 2 ; dois, com conceito 3 ; dois, com conceito 4 ; e cinco ficaram com conceito 5.
Nenhum ficou sem conceito ( SC ) nessa região.
Também na região Centro-Oeste, a maioria dos cursos foi \textbf{ofertada} na modalidade presencial : 46 ( 97,9 \% ).
Para esses, a moda foi o conceito 3, com 21 cursos ( 45,7 \% ).
Os demais cursos presenciais da região foram alocados nos conceitos 2 ( 13 cursos, 28,3 \% ), 5 ( cinco cursos, 10,9 \% ), 4 ( quatro cursos, 8,7 \% ) e 1 ( dois cursos, 4,3 \% ).
Um curso presencial da região ficou sem conceito ( SC ).
Na modalidade a distância, foi oferecido um curso com conceito 2.
Nenhum curso recebeu conceitos 1, 3, 4 ou 5.
Nenhum curso ficou sem conceito ( SC ) nessa região.
Para o único curso a distância nessa região, o conceito atribuído foi 3.
A título de síntese, no Gráfico 4.2, apresenta -se a distribuição do Conceito Enade, segundo a categoria administrativa da IES.
Nesse gráfico, como dito anteriormente, \textbf{poligonais} mais à esquerda representam Conceito Enade mais baixos, e \textbf{poligonais} mais à direita, Conceito Enade mais altos.
É importante observar que, para cada \textbf{poligonal} presente nesse gráfico, a porcentagem constante no eixo y refere -se ao universo de cursos oferecidos por instituições públicas ( \textbf{poligonal} azul ) ou privadas ( \textbf{poligonal} verde ), respectivamente.
Os cursos em IES públicas ( \textbf{poligonal} azul ) apresentam uma \textbf{poligonal} mais à direita do que os em IES privadas ( \textbf{poligonal} verde ) e, consequentemente, uma distribuição de Conceito Enade com valores maiores. 4.3 CONCEITOS POR ORGANIZAÇÃO ACADÊMICA E GRANDE REGIÃO Na Tabela 4.3, encontra -se a distribuição dos conceitos atribuídos aos cursos participantes do Enade 2023, na área de Arquitetura e Urbanismo, por organização acadêmica, segundo as grandes regiões brasileiras.
Como mostram os dados da Tabela 4.3, dos 564 cursos de Arquitetura e Urbanismo participantes, 208 eram oferecidos em universidades, 223 em centros universitários e 123 em faculdades ; dez cursos foram oferecidos por CEFET/IF.
Essa distribuição corresponde a, respectivamente, 36,9 \%, 39,5 \%, 21,8 \% e 1,8 \% dos cursos.
De acordo com os dados apresentados na mesma tabela, de todos os 97 cursos avaliados com o conceito 4, 53 eram vinculados a universidades.
Os cursos vinculados a esse tipo de organização acadêmica tiveram o conceito 3 como conceito modal, atribuído a 82 cursos ( 39,4 \% ).
Os demais cursos \textbf{ofertados} por universidades receberam os conceitos 2 ( 39 cursos ), 5 ( 31 cursos ) e 1 ( um curso ).
Dois cursos ficaram sem conceito ( SC ).
Dos cursos vinculados a centros universitários, o conceito modal também foi o 3, atribuído a 121 cursos ( 54,3 \% ).
Seis cursos vinculados a esse tipo de organização acadêmica ficaram sem conceito ( SC ), e os outros cursos receberam os conceitos 2 ( 58 cursos ), 4 ( 30 cursos ), 1 ( seis cursos ) e 5 ( dois cursos ).
Dos 123 cursos mantidos por faculdades, 52 ( 42,3 \% ) receberam conceito 3, o conceito modal.
O segundo conceito mais frequentemente atribuído aos cursos oferecidos por essa categoria de IES foi o conceito 2, recebido por 48 cursos.
Na sequência, dez cursos obtiveram conceito 4 ; seis cursos, o conceito 1 ; e dois cursos, o conceito 5.
Dos cursos oferecidos em faculdades, cinco cursos ficaram sem conceito ( SC ).
Dos cursos oferecidos por CEFET/IF, o conceito modal foi o 4, atribuído a quatro cursos ( 40,0 \% ).
Nenhum curso vinculado a esse tipo de organização acadêmica ficou sem conceito ( SC ), e os outros cursos receberam os conceitos 3 ( quatro cursos ), 5 ( um curso ), 2 ( um curso ) e 1 ( nenhum curso ).
Considerando -se, separadamente, as regiões brasileiras, verifica -se que, na região Norte, as universidades participaram com 11 cursos avaliados, aos quais foram atribuídos os seguintes conceitos : 1 a nenhum curso ; 2 a quatro cursos ; 3, o conceito modal, a quatro cursos ; conceito 4 a três cursos ; e 5 a nenhum curso.
Nenhum curso ficou sem conceito ( SC ).
Os centros universitários da região Norte foram representados por 18 cursos, dos quais um curso obteve conceito 1 ; sete cursos, conceito 2 ; nove cursos, conceito 3, o valor modal ; um curso, o conceito 4 ; e nenhum curso, o conceito 5.
Nenhum curso ficou sem conceito ( SC ).
As faculdades participaram com dez cursos na região Norte, os quais receberam os conceitos 1 ( um curso ), 2 ( três cursos ), 3 ( cinco cursos ), 4 ( nenhum curso ) e 5 ( nenhum curso ).
Um curso ficou sem conceito ( SC ).
Os CEFET/IF participaram com um curso, ao qual foi atribuído o conceitos 3.
Na região Nordeste, as universidades participaram com 23 dos 95 cursos da área de Arquitetura e Urbanismo oferecidos.
O conceito modal foi 3, atribuído a nove cursos.
Outros sete cursos obtiveram o conceito 4 ; seis cursos, o conceito 5 ; e nenhum curso, o conceito 1.
Nenhum curso ficou sem conceito ( SC ).
Os centros universitários \textbf{contaram} com 44 cursos participantes na região Nordeste, 18 dos quais receberam o conceito modal 3.
Para os demais cursos, foram atribuídos os conceitos 2 ( 15 cursos ), 4 ( seis cursos ), 1 ( dois cursos ) e 5 ( um curso ).
Dois cursos oferecidos por centros universitários no Nordeste ficaram sem conceito ( SC ).
As faculdades foram representadas por 26 cursos na região Nordeste, os quais receberam os conceitos 1 ( um curso ), 2 ( nove cursos ), 3 ( 13 cursos ), 4 ( dois cursos ) e 5 ( nenhum curso ).
Um curso ficou sem conceito ( SC ).
Na região Nordeste, os CEFET/IF participaram com dois cursos, aos quais foi atribuído o conceito 3.
Na região Sudeste, as universidades concentraram 91 dos 246 cursos de Arquitetura e Urbanismo da região.
Dos cursos oferecidos em IES com este tipo de organização acadêmica, nessa região, o conceito modal foi 3, atribuído a 38 cursos.
Os demais cursos receberam os conceitos 4 ( 21 cursos ), 2 ( 19 cursos ), 5 ( 12 cursos ) e 1 ( nenhum curso ).
Um curso ficou sem conceito ( SC ).
Os centros universitários participaram com 96 cursos na região Sudeste, dos quais 54 obtiveram o conceito modal 3.
O conceito 2 foi atribuído a 24 cursos ; o conceito 4 a 14 cursos ; o conceito 1 a dois cursos ; e o conceito 5 a nenhum curso.
Dois cursos ficaram sem conceito ( SC ) nessa combinação de organização acadêmica e região.
Já as faculdades foram representadas por 55 cursos na região Sudeste, dos quais dois ficaram sem conceito ( SC ).
O conceito modal 2 foi atribuído a 24 cursos.
Aos demais cursos foram atribuídos os conceitos 3 ( 19 cursos ), 4 ( seis cursos ), 5 ( dois cursos ) e 1 ( dois cursos ).
Os CEFET/IF participaram com quatro cursos na região Sudeste, aos quais foram atribuídos os conceitos 4 a três cursos e 5 a um curso.
Nenhum curso ficou sem conceito ( SC ).
Dos 136 cursos da região Sul, 66 eram mantidos por universidades, para os quais o conceito modal foi 3, atribuído a 25 cursos.
Os demais cursos receberam os conceitos 4 ( 19 cursos ), 2 ( 12 cursos ), 5 ( oito cursos ) e 1 ( um curso ).
Em IES com esse tipo de organização acadêmica, um curso da região Sul ficou sem conceito ( SC ).
Dos 47 cursos participantes de centros universitários da região Sul, um ficou sem conceito ( SC ), e 28 cursos receberam conceito 3, o valor modal.
Os conceitos \textbf{restantes} ficaram assim distribuídos : 4 ( oito cursos ), 2 ( oito cursos ), 5 ( um curso ) e 1 ( um curso ).
Quanto aos 21 cursos vinculados a faculdades na região Sul, 13 cursos receberam o conceito 3, valor modal.
Os demais conceitos ficaram assim distribuídos : 2 ( cinco cursos ), 4 ( dois cursos ), 5 ( nenhum curso ) e 1 ( nenhum curso ).
Dos 21 cursos, um ficou sem conceito ( SC ).
Os CEFET/IF participaram com dois cursos na região Sul, aos quais foram atribuídos os seguintes conceitos : 1 a nenhum curso ; 2 a um curso ; 3 a nenhum curso ; 4 a um curso ; e 5 a nenhum curso.
Nenhum curso ficou sem conceito ( SC ).
Na região Centro-Oeste, 17 dos 47 cursos oferecidos eram mantidos por universidades.
Para esse tipo de organização acadêmica, o conceito modal foi 3, atribuído a seis cursos.
Em relação às universidades, nenhum curso da região ficou sem conceito ( SC ).
Os demais cursos receberam os conceitos 5 ( cinco cursos ), 4 ( três cursos ), 2 ( três cursos ) e 1 ( nenhum curso ).
Os centros universitários da região Centro-Oeste \textbf{contaram} com 18 cursos : 12 cursos foram avaliados com o conceito 3 ( modal ) ; quatro cursos, com o conceito 2 ; um curso, com o conceito 4 ; nenhum curso, com o conceito 5 ; e nenhum curso obteve o conceito 1.
Um curso ficou sem conceito ( SC ).
Dos 11 cursos oferecidos em faculdades na região Centro-Oeste, sete cursos receberam o conceito 2 ( modal ).
Os conceitos \textbf{restantes} ficaram assim distribuídos : 3 ( dois cursos ), 1 ( dois cursos ), 5 ( nenhum curso ) e 4 ( nenhum curso ).
Nessa região, nenhum curso oferecido em faculdades ficou sem conceito ( SC ) ou obteve os conceitos 1 e 5.
No que se refere aos CEFET/IF da região Centro-Oeste, um curso foi avaliado com o conceito 3 ( modal ).
Nenhum curso recebeu os demais conceitos ou ficou sem conceito ( SC ).
De maneira semelhante à forma como se \textbf{procedeu} em relação à seção anterior, esta seção encerra -se com o Gráfico 4.3, que apresenta a distribuição do Conceito Enade dos cursos de Arquitetura e Urbanismo segundo a organização acadêmica da IES.
No Gráfico 4.3, apresenta -se a distribuição do Conceito Enade segundo a organização acadêmica da IES.
A leitura desse último gráfico é semelhante aos anteriores, seguindo -se a lógica de que \textbf{poligonais} deslocadas para a esquerda correspondem aos conceitos mais baixos ( 1, 2 e 3 ) e \textbf{poligonais} deslocadas para a direita correspondem a Conceitos Enade mais altos.
Nas universidades ( linha azul ), todos os cursos obtiveram conceito 5 ou abaixo disso, já que a \textbf{poligonal} \textbf{cumulativa} alcança 100 \% para o conceito 5.
Além disso, os cursos de CEFET/IES apresentam uma \textbf{poligonal} mais à direita do que as demais, \textbf{denotando} uma melhor distribuição de Conceitos Enade.
No outro extremo, encontram -se os cursos mantidos por faculdades ( linha vermelha ), apresentando as respectivas \textbf{poligonais} mais à esquerda e, consequentemente, uma distribuição de Conceitos Enade com os menores valores.
Os cursos mantidos por universidades ( linha azul ) e centros universitários ( linha laranja ) apresentam \textbf{poligonais} mais centrais e, consequentemente, uma distribuição de Conceitos Enade com os valores intermediários.

% matched lemmas: abscissa, contar, cruzar, cumulativo, denotar, equivaler, finalidade, indicado, ministrar, ofertar, ordenado, panorama, poligonal, pontilhado, proceder, reiterativo, restante, roxa, segunda, tabular, tipologia, variável, visualizar
\end{textsample}
